\section{Introduction}
\label{sec:intro}

Ocean warming is one of the most pervasive stressors acting on marine ecosystems \citep{hobday2016}.
In addition to gradual warming, discrete marine heatwaves have increased in frequency and duration over the last century \citep{oliver2018}.
Copepods are the most abundant metazoans in the ocean and they play a key role in the transfer of energy from primary producers to higher trophic levels \citep{turner2004}.

Despite this, it is not well known how copepods respond when they experience several heatwaves one after another.
Previous work has mostly focused on single exposures \citep{missingref2020}, which may underestimate cumulative effects.
In this study we ask whether repeated exposure leads to acclimation or to accumulated damage.
